% Generated by analysis/tables/make_all_tables.py -- do not edit by hand.
\begin{table}[H]
\centering
\small
\caption{Mechanism placebo comparison}
\label{tab:mech_placebos}
\begin{threeparttable}
\begin{tabular}{lrrr}
\toprule
Metric & Real shutdown & Mean placebo & Maximum placebo \\
\midrule
Released in next 6 days & 25.38\% & 0.56\% & 0.77\% \\
Delay $>$24h & 29.71\% & 1.10\% & 1.31\% \\
Delay $>$48h & 25.38\% & 0.72\% & 0.90\% \\
Delay $>$72h & 22.75\% & 0.46\% & 0.56\% \\
Delay $>$120h & 13.88\% & 0.32\% & 0.46\% \\
\bottomrule
\end{tabular}
\begin{tablenotes}[flushleft]
\footnotesize\raggedright
\item \textit{Notes:} Placebo windows are six six-day authoring windows between July 1 and July 11, with six-day subsequent release windows that end before treatment. Windows overlap, so the comparison is descriptive rather than an exact independent randomization test. Commit metadata were queried for all first-push commits pushed on or after July 18 but for a daily sample of those pushed July 1--17 (7--17 percent of commits per push date). Placebo rates therefore weight each commit by the inverse of its push date's query rate (number of first-push commits divided by number queried); commits authored July 18--23 can only be pushed on or after July 18 and receive weight one, so all rows estimate the same population share. Without weights, placebo rates would over-represent late pushes. Mean placebo $N$ is 4,144 observed commits per window. Sources: \texttt{results/final/table\_mechanism\_placebos.csv}, reproduced by \texttt{analysis/mechanism/01\_placebo\_weighting\_audit.py}.
\end{tablenotes}
\end{threeparttable}
\end{table}
